\documentclass{article}
\usepackage{iclr2027_conference,times}

\input{math_commands.tex}

\usepackage{amsmath,amssymb,booktabs,microtype,tabularx,xcolor}
\usepackage[hidelinks]{hyperref}
\usepackage{url}

\title{Beyond One-Shot Attribution: Multi-Round Bilevel Data Selection}

\author{Anonymous Authors}

\newcommand{\Dtr}{\mathcal{D}_{\mathrm{tr}}}
\newcommand{\Dval}{\mathcal{D}_{\mathrm{val}}}
\newcommand{\Wbudget}{\mathcal{W}_{B}}
\newcommand{\thetahat}{\widehat{\vtheta}}
\newcommand{\thetatilde}{\widetilde{\vtheta}}
\newcommand{\thetaIL}{\vtheta_{\mathrm{IL}}}
\newcommand{\rhoideal}{s_{\mathrm{RHO}}^{\mathrm{ideal}}}
\newcommand{\rhopractical}{s_{\mathrm{RHO}}^{\mathrm{practical}}}

\begin{document}

\maketitle

\begin{abstract}
Data attribution evaluates a local change in training-example utility, whereas
data selection must optimize a feasible subset along a changing training
trajectory. We connect the two through a bilevel program whose outer state is
a continuous, fixed-budget weight vector. The value-function penalty used by
fully first-order stochastic bilevel optimization (F2SA) yields a two-model
per-example loss difference; idealized Reducible Holdout Loss Selection
(RHO-LOSS) gives the same direction at a fixed state, while practical RHO-LOSS
also changes the comparator, candidate scope, scoring frequency, and training
timing. This view turns ``attribution then selection'' into a controlled
multi-round question. Under a fixed target-training horizon, we compare one to
five complete inner-solve/outer-update rounds, reset and persistent continuous
weights, and matched random controls. We then factor practical RHO into global
versus batch-local quotas, stale versus online scores, immediate updates, and
frozen versus updating comparators, including a dedicated batch-size scaling
study. The resulting design tests whether gains come from optimizing utility
beyond the uniform point, from merely recomputing a local score, or from the
online protocol rather than the score formula itself.
\end{abstract}

\section{Introduction}

Data attribution and data selection are often evaluated with the same
rank-and-retrain pipeline, but they answer different questions. Attribution
asks how an objective changes under a local perturbation of one training
example. Selection asks which budget-feasible data should shape an entire
training trajectory. A score evaluated once at uniform weights can therefore
be a useful direction without being a solution to the selection problem.

We make this distinction explicit through bilevel optimization over continuous
example weights. In this coordinate system, the value-function construction of
F2SA \citep{kwon2023f2sa} produces the difference between losses under two
lower-level solutions. Idealized RHO-LOSS \citep{mindermann2022rholoss}
produces the same per-example direction at a fixed outer state. Practical
RHO-LOSS is farther away: it replaces the updating joint solution with a
frozen holdout model, scores a batch-local candidate pool, takes a local quota,
and immediately updates a nonconverged target. These changes can be useful,
but their effects should not be attributed to the loss-difference formula
without a controlled decomposition.

The empirical question is thus not simply whether a value-function score beats
RHO-LOSS. It is whether optimizing the same continuous outer state beyond the
uniform point produces better fixed-budget data, and which parts of the
practical online protocol explain any remaining advantage. We answer this with
a fixed-horizon block design. After a shared warm start, every target runs for
five equal blocks. A selector may update at the first one through five block
boundaries, but all branches train for the same total epochs and examples. The
design directly compares one-shot selection, repeated scores reset to uniform,
and persistent projected optimization. It therefore distinguishes outer
memory from the simpler benefit of rescoring a changed model.

Our contributions are:
\begin{enumerate}
    \item We derive, in one native notation, the exact correspondence between
    the F2SA value-function direction and idealized RHO-LOSS, while marking the
    assumptions broken by the practical frozen-model estimator.
    \item We formulate one-shot attribution and multi-round selection as
    $R=1$ and $R=2,\ldots,5$ instances of one fixed-budget, fixed-horizon
    protocol, with reset and persistent outer states and matched random
    controls.
    \item We separate RHO's score approximation from its training protocol by
    adding global scoring, batch-local quota, immediate update, and online
    rescoring one at a time, then measure how the online batch scale changes
    accuracy, optimizer steps, coverage, and computation.
\end{enumerate}
The paper's main claim is deliberately conditional: persistent multi-round
optimization is supported only if it improves over both its one-shot branch
and the corresponding reset branch at equal target-training horizon. If a
batch-faithful RHO variant wins instead, the decomposition determines whether
the source is the frozen comparator, local diversity constraint, online
adaptation, or optimizer-step granularity.

\section{Budgeted Data Selection as Bilevel Optimization}
\label{sec:formulation}

Let $\Dtr=\{z_i\}_{i=1}^{n}$ be candidate training examples and let
$\Dval$ be a clean holdout set. We use
\begin{align}
g(\vomega,\vtheta)
    &:= \frac{1}{n}\sum_{i=1}^{n}\omega_i\ell(\vtheta;z_i),
    &
f(\vtheta)
    &:= \frac{1}{|\Dval|}\sum_{z\in\Dval}\ell(\vtheta;z),
\label{eq:fg}
\end{align}
and define
\begin{equation}
\min_{\vomega\in\Wbudget} h(\vomega)
    := f(\thetahat(\vomega))
\quad\text{subject to}\quad
\thetahat(\vomega)\in
    \argmin_{\vtheta}g(\vomega,\vtheta),
\label{eq:bilevel-selection}
\end{equation}
where
\begin{equation}
\Wbudget:=\left\{\vomega\in[0,1]^n:
\sum_{i=1}^{n}\omega_i=B\right\}.
\label{eq:capped-simplex}
\end{equation}
The continuous vector $\vomega$ is the outer optimization state. The target
model is trained on the discrete set
$S(\vomega)=\operatorname{TopB}(\vomega)$ only when a training decision is
required. The binary mask is not substituted back for $\vomega$ in the next
outer update.

The budget-feasible uniform point is
\begin{equation}
\vomega^{\mathrm{unif}}=\frac{B}{n}\bm{1}.
\label{eq:uniform-budget}
\end{equation}
It is the appropriate reference for comparing one-shot and iterative
selection. The all-ones vector has a useful attribution interpretation---all
examples are fully present---but lies outside $\Wbudget$ whenever $B<n$.
Conflating these two reference points creates an artificial first projection
that can dominate the score update.

Equation~\ref{eq:bilevel-selection} is our optimization target, not a claim
that the nonconvex neural-network problem satisfies the assumptions used in
all bilevel convergence analyses. In particular, the F2SA theory assumes an
unconstrained outer variable and a strongly convex lower problem. We use its
value-function construction to derive an estimator, and treat projection onto
the set in Eq.~(\ref{eq:capped-simplex}) as an explicit algorithmic adaptation whose behavior
must be tested empirically.

F2SA's MNIST hyper-cleaning experiment instead optimizes unconstrained logits
$\va\in\mathbb{R}^n$ and uses sigmoid weights $\omega_i=\sigma(a_i)$; it does
not impose a fixed cardinality budget. To separate parameterization from the
estimator, our primary method optimizes $\vomega$ directly with capped-simplex
projection. As an ablation, we use
\begin{equation}
\omega_i=\sigma(a_i+c(\va)),
\qquad \sum_i\sigma(a_i+c(\va))=B,
\end{equation}
where the scalar shift $c(\va)$ is obtained by one-dimensional root finding.
Learning rates are tuned separately because direct weights and logits have
different geometries.

\section{A Value-Function View of RHO-LOSS}
\label{sec:connection}

\subsection{The F2SA value-function direction}

Define the lower-level value function
\begin{equation}
g^\star(\vomega):=\min_{\vtheta}g(\vomega,\vtheta)
=g(\vomega,\thetahat(\vomega)),
\qquad
\thetahat(\vomega)\in\argmin_{\vtheta}g(\vomega,\vtheta).
\end{equation}
The bilevel feasibility condition is
$g(\vomega,\vtheta)-g^\star(\vomega)\leq 0$. We therefore consider the
penalized value-function objective from the original formulation,
\begin{equation}
\mathcal{L}^{\alpha}(\vomega,\vtheta,\vtheta')
=f(\vtheta)+\alpha\bigl(g(\vomega,\vtheta)
-g(\vomega,\vtheta')\bigr),
\qquad \alpha>0,
\label{eq:f2sa-penalty}
\end{equation}
and track two solutions,
\begin{align}
\thetatilde(\vomega)
&\in\argmin_{\vtheta}
    \left[f(\vtheta)+\alpha g(\vomega,\vtheta)\right],
\label{eq:tilde-solution}\\
\thetahat(\vomega)
&\in\argmin_{\vtheta}g(\vomega,\vtheta).
\label{eq:hat-solution}
\end{align}
Under the regularity conditions needed for the envelope theorem, the gradient
of the minimized penalty is
\begin{equation}
\nabla_{\vomega}\mathcal{L}^{\alpha,\star}(\vomega)
=\alpha\left[
\nabla_{\vomega}g(\vomega,\thetatilde)
-\nabla_{\vomega}g(\vomega,\thetahat)
\right].
\label{eq:f2sa-gradient}
\end{equation}
For the linear weighting in Eq.~\eqref{eq:fg}, its $i$th coordinate is
\begin{equation}
\left[\nabla_{\vomega}\mathcal{L}^{\alpha,\star}(\vomega)\right]_i
=\frac{\alpha}{n}\left[
\ell(\thetatilde;z_i)-
\ell(\thetahat;z_i)
\right].
\label{eq:loss-difference-gradient}
\end{equation}
Consequently, the loss difference
\begin{equation}
s_i^{\mathrm{VF}}(\vomega)
:=\ell(\thetahat(\vomega);z_i)
-\ell(\thetatilde(\vomega);z_i)
\label{eq:vf-score}
\end{equation}
is proportional to the \emph{negative} penalized outer gradient. A selector
that keeps examples with large $s_i^{\mathrm{VF}}$ therefore follows a descent
direction, subject to inner-solution and finite-$\alpha$ errors. The same
$\alpha$ must appear in both the penalized inner objective and the outer
gradient. A fixed value is a finite-penalty approximation; a scheduled value
must be coordinated with the inner and outer step sizes. Score
standardization removes only the global scale of Eq.~\eqref{eq:f2sa-gradient},
not the dependence of $\thetatilde$ on $\alpha$.

\subsection{Ideal and practical RHO-LOSS}

RHO-LOSS begins with online batch selection. At training step $t$, a large
candidate batch $C_t$ is sampled, a score is computed for every example in
$C_t$, and only the top $b$ examples are used for one target-model update. Its
Bayesian derivation targets the example that would most reduce predictive
holdout loss. After replacing Bayesian posteriors with trained point-estimate
models, the idealized score can be written in our notation as
\begin{equation}
\rhoideal_i(\vomega)
=\ell(\thetahat(\vomega);z_i)
-\ell(\thetatilde(\vomega);z_i),
\label{eq:ideal-rho}
\end{equation}
where the second model incorporates both the holdout objective and the data
represented by $\vomega$. Comparing Eqs.~\eqref{eq:vf-score} and
\eqref{eq:ideal-rho} gives
\begin{equation}
\rhoideal(\vomega)
=s^{\mathrm{VF}}(\vomega)
=-\frac{n}{\alpha}
\nabla_{\vomega}\mathcal{L}^{\alpha,\star}(\vomega).
\label{eq:rho-vf-equivalence}
\end{equation}
This is an equality of directions at a fixed $\vomega$ under exact inner
solutions and aligned objective normalizations. It does not say that one
Top-$b$ RHO step solves Eq.~\eqref{eq:bilevel-selection}.

Practical RHO-LOSS makes an additional approximation. It trains an
irreducible-loss model once on $\Dval$,
\begin{equation}
\thetaIL\in\argmin_{\vtheta}f(\vtheta),
\end{equation}
freezes it, and uses
\begin{equation}
\rhopractical_i(t)
=\ell(\vtheta_t;z_i)-\ell(\thetaIL;z_i),
\label{eq:practical-rho}
\end{equation}
where $\vtheta_t$ is the current, generally nonconverged target model. This
removes the need to update the joint comparator as selected data change. It
also means Eq.~\eqref{eq:practical-rho} is not, in general, the exact gradient
in Eq.~\eqref{eq:rho-vf-equivalence}. The RHO-LOSS paper finds this frozen
approximation computationally useful and sometimes empirically preferable;
our framework treats that as a testable approximation rather than an
identity.

\subsection{The approximation ladder}

The shared formulation yields the following ladder:
\begin{equation}
\underbrace{\nabla h(\vomega)}_{\text{bilevel target}}
\;\longleftarrow\;
\underbrace{\nabla\mathcal{L}^{\alpha,\star}(\vomega)}_{\text{finite-penalty VF}}
\;\longleftrightarrow\;
\underbrace{-\rhoideal(\vomega)}_{\text{two exact inner models}}
\;\longrightarrow\;
\underbrace{-\rhopractical(t)}_{\text{frozen IL proxy}}.
\label{eq:approximation-ladder}
\end{equation}
The left arrow carries finite-penalty bias and relies on F2SA's regularity and
tracking conditions. The middle equivalence follows algebraically from
linear example weights. The right arrow introduces the frozen comparator and
target-trajectory approximations. Batch-local ranking and Top-$B$
discretization are further approximations applied after the score is formed.

\section{From a Local Attribution Direction to Multi-Round Selection}
\label{sec:single-to-multi}

Equation~\ref{eq:rho-vf-equivalence} is a local statement: at a fixed
$\vomega$, a two-model loss difference supplies a descent direction for a
finite-penalty surrogate. Data selection is instead a constrained, path-
dependent decision problem. It must specify an initial state, how long the
real target is trained before the selector changes, whether the same outer
state is retained, and how a continuous state becomes a discrete subset.

\paragraph{A common warm start.}
All comparisons begin with the same target initialization and the same random
budget-$B$ subset $S_0$. The target trains on $S_0$ for $E_0$ epochs before any
method-dependent selection. This avoids choosing an arbitrary Top-$B$ set from
tied uniform weights and makes the first comparison occur at an identical,
nontrivial target state.

\paragraph{One fixed horizon, one to five selection rounds.}
After the warm start, target training is partitioned into five equal blocks of
$L$ epochs, for a common horizon $H=5L$. At the start of a block, one complete
selection round (i) updates $\thetahat$ and $\thetatilde$ under a fixed
inner-pass rule, (ii) computes a global direction, (iii) updates and projects
the continuous state, and (iv) deploys
$S_{r}=\operatorname{TopB}(\vomega_r)$ for real target training. We compare
$R\in\{1,2,3,4,5\}$: the first $R$ boundaries update the selector, after which
its continuous state and subset are frozen while the target still trains to
$H$. An extra round is therefore a complete bilevel selection update, not an
extra target-SGD step or a virtual epoch that is discarded.

\paragraph{Reset separates rescoring from outer memory.}
A reset branch computes every new direction from
$\vomega^{\mathrm{unif}}$; a persistent branch updates the same
$\vomega_r\in\Wbudget$ across boundaries. If both improve with $R$, repeated
evaluation at a changing learner state is sufficient. If only the persistent
branch improves, the evidence instead favors optimization of continuous
weights along a path. In both cases Top-$B$ is a deployment operator: the
binary mask used to train the target is never used as the next differentiable
outer state.

\paragraph{Where can a multi-round gain come from?}
Relative to one-shot selection, four state changes can matter:
\begin{enumerate}
    \item the target model changes, so the value of an example is conditional
    on a new learner state;
    \item the two inner models track new finite-penalty solutions;
    \item the continuous outer state remembers previous directions; and
    \item Top-$B$ is applied again, changing the deployed discrete subset.
\end{enumerate}
The fixed-horizon reset/persistent comparison identifies the first three at
the method level; score, state, subset, and target checkpoints at every block
boundary expose the fourth. A claim that ``more rounds help'' is incomplete
unless the responsible transition is identified.

\paragraph{RHO as an approximation endpoint.}
Practical RHO removes two costly pieces of the multi-round VF procedure: it
uses the current nonconverged target instead of $\thetahat(\vomega)$ and a
frozen holdout-only IL model instead of
$\thetatilde(\vomega)$. It then scores only a random large batch and
immediately trains on its top fraction. We therefore compare RHO neither as an
unrelated baseline nor as an exact VF implementation, but as the efficient end
of a nested path. Global scoring, local quota, immediate updates, online
rescoring, and comparator choice are switched on separately in
Section~\ref{sec:experiments}.

\section{Research Questions}
\label{sec:rqs}

\paragraph{RQ1: Uniform utility versus optimized utility.}
Does replacing one attribution update at
$\vomega^{\mathrm{unif}}$ with a persistent sequence of continuous,
budget-feasible weights produce a better Top-$10\%$ subset? We answer with the
paired gap between persistent and reset branches at each common freeze point.
An improvement over one-shot alone is insufficient: persistent optimization
must also beat repeated reset scoring to establish a benefit from outer
memory. The supporting diagnostics are development loss, selected-noise rate,
weight movement, subset turnover, and final performance at the common horizon.

\paragraph{RQ2: How does selection scale with complete bilevel rounds?}
When the full inner-update, outer-score, weight-projection, and Top-$B$
deployment cycle is repeated $R\in\{1,2,3,4,5\}$ times, does performance
improve monotonically or saturate? Every branch trains the target to the same
$H=5L$ horizon, so $R$ changes selection opportunities but not real target
epochs or examples. We report the entire freeze-$R$ curve rather than choosing
one round count on test data, and pair accuracy-versus-target-data with
accuracy-versus-total-compute.

\paragraph{RQ3: Which technical choices create practical RHO's gains?}
We factor RHO into four protocol choices---global versus batch-local quota,
blockwise versus immediate target updates, stale versus online scores, and an
updating joint comparator versus a frozen holdout IL model. At matched protocol
points we also substitute the strict VF, updating-comparator, and frozen-IL
score formulas. This design asks separately whether a gain comes from a more
useful loss difference, greater batch diversity, faster adaptation, or finer
optimizer-step granularity. Batch-size scaling is the priority experiment
because it simultaneously controls the size of the local ranking pool and the
number of immediate target updates.

The algebraic correspondence in Section~\ref{sec:connection} is treated as a
correctness precondition, not a standalone performance RQ. Gradient sign,
inner tracking, budget projection, and ranking agreement must pass the gates in
Section~\ref{sec:experiments} before the three questions above receive a
value-function interpretation.

\section{Experimental Design}
\label{sec:experiments}

\subsection{Stage 0: correctness gates}

Before training image classifiers, we use a low-dimensional strongly convex
problem for which the exact hypergradient is available. We compare (i) the
analytic hypergradient, (ii) finite differences, (iii) fully unrolled
autodifferentiation, and (iv) Eq.~\eqref{eq:loss-difference-gradient}. The
sign must agree and pairwise cosine similarity must exceed $0.99$. We also
test that projection preserves $0\leq\omega_i\leq1$ and
$\sum_i\omega_i=B$, that the uniform state equals $B/n$, and that tie breaking
is deterministic.

On MNIST with 10\% label noise, we then hold $\vomega$ fixed and increase the
number of inner updates. We record both inner optimality residuals, score
cosine, Top-$B$ Jaccard similarity, and agreement with an independently
retrained solution. The main study starts only after the selected stopping
rule produces a stable score. Validation data determine stopping and
hyperparameters; the test set is never used for either.

\subsection{RQ1: correspondence and approximation diagnostics}

At matched checkpoints and the same $\vomega$, we construct the following
scores:
\begin{enumerate}
    \item exact or fully unrolled hypergradient on the tractable problem;
    \item converged finite-penalty VF score, Eq.~\eqref{eq:vf-score};
    \item truncated/warm-start VF score after a controlled number of inner
    updates;
    \item ideal RHO score with an updating joint comparator;
    \item practical RHO score with a frozen holdout-only IL model.
\end{enumerate}
We report signed cosine similarity, Spearman rank correlation, Top-$B$
overlap, and selected-noise fraction. We sweep the penalty multiplier
$\alpha$, or a preregistered schedule $\alpha_r$, and report the associated
inner residuals. This separates finite-penalty bias from optimization and
frozen-comparator errors.

\subsection{RQ2: one-shot versus multi-round optimization}

All methods begin from the same model initialization and the same random
budget-$B$ subset $S_0$. After the common warm start, we compare:
\begin{center}
\begin{tabular}{lll}
\toprule
ID & Method & Purpose \\
\midrule
U-fixed & fixed random subset & no reselection baseline \\
U-round & random subset each round & reselection-only baseline \\
RHO-1 & one practical RHO selection & one-shot practical attribution \\
VF-1 & one verified VF step & one-shot value-function direction \\
RHO-M & repeated practical RHO & practical multi-round selection \\
VF-M & persistent projected VF & iterative outer optimization \\
\bottomrule
\end{tabular}
\end{center}
The outer update interval is the number of real target epochs between score
updates,
\begin{equation}
\tau\in\{1,5,10,20,\infty\},
\end{equation}
where $\tau=\infty$ performs one selection only after the prescribed target
training block. This variable replaces the earlier ``persistence horizon''
and ``online steps'' sweeps: it directly controls when a new selector state can
affect the real training trajectory. At each finite $\tau$, the continuous
state persists; only a separate reset ablation returns it to
$\vomega^{\mathrm{unif}}$.

The primary estimand is the paired multi-round gain
\begin{equation}
\Delta_{\mathrm{multi}}
=A(\mathrm{VF\mbox{-}M})-A(\mathrm{VF\mbox{-}1}),
\end{equation}
at equal target data and optimizer-step budgets. We also report target
accuracy versus selected examples, accuracy versus wall-clock/compute,
selected-noise fraction, class coverage, and selection turnover.

\subsection{RQ3: nested approximation ablations}

Starting from verified VF-M, we introduce one approximation at a time:
\begin{enumerate}
    \item reduce inner updates while keeping $\alpha$ and $\vomega$ fixed;
    \item freeze the joint comparator, then replace it with a holdout-only IL
    model;
    \item replace global scoring with RHO's batch-local $n_B\!\to n_b$
    protocol;
    \item delay Top-$B$ discretization while retaining continuous updates;
    \item reset the continuous state after each selection round.
\end{enumerate}
We additionally compare capped-simplex projection with the budget-calibrated
sigmoid parameterization from Section~\ref{sec:formulation}; this is a
parameterization ablation, not a different selection budget.
Four diagnostic controls isolate the multi-round mechanism: frozen target
state, frozen inner states, reset $\vomega$, and delayed Top-$B$. Their paired
differences decompose gains due to changes in the target trajectory, inner
tracking, persistent outer optimization, and repeated discretization.

Target and inner batch sizes are separate variables. We first screen
\texttt{target\_batch\_size}$\in\{32,128,\text{full selected set}\}$ with a
fixed inner batch, then screen
\texttt{inner\_batch\_size}$\in\{32,128,320,\text{full candidate set}\}$ at the
chosen target batch. Comparisons match examples processed and separately
report optimizer steps. Only the two best configurations are repeated with
all seeds.

The faithful practical RHO baseline follows the original online protocol:
sample a large batch of $n_B=320$, compute frozen-IL reducible losses, select
$n_b=32$, and immediately update the target on those examples. Selection and
target-update BatchNorm statistics are computed on the large and selected
batches, respectively. Global RHO scoring is reported as our extension rather
than as the original RHO-LOSS algorithm.

\subsection{Datasets, budgets, and staged execution}

We proceed from MNIST to CIFAR-10 and CIFAR-100. The initial study uses clean
and 10\% symmetric-label-noise conditions, 10\% retention, and three paired
seeds. We expand to 20\% and 30\% noise only after the direction is stable
across at least two of three seeds. After fixing the 10\% protocol, we evaluate
retention in $\{5\%,10\%,15\%,20\%\}$ for U-round, VF-1, VF-M, and the best
practical RHO variant.

Every outer update stores the continuous weights, selected indices, raw and
standardized scores, weight sum/minimum/maximum/effective sample size, both
inner objectives and residuals, score cosine, Top-$B$ Jaccard similarity,
class histogram, selected-noise fraction, target metrics, processed examples,
optimizer steps, wall-clock time, and peak memory. A run is admissible only
if its configuration, seeds, initial subset $S_0$, and all state transitions
are reproducible.

\section{Related Work}

\paragraph{Attribution and bilevel example reweighting.}
Training-data attribution assigns examples credit for a prediction or utility:
influence functions differentiate an infinitesimal reweighting at a fitted
model \citep{koh2017influence}, TracIn accumulates first-order influence along
saved training checkpoints \citep{pruthi2020tracin}, and Data Shapley averages
marginal utility over data coalitions \citep{ghorbani2019datashapley}. A
complementary line directly learns training weights from clean validation
loss, either through a one-step meta-gradient \citep{ren2018reweight} or by
unrolling an approximate bilevel problem \citep{franceschi2018bilevel}.
F2SA removes Hessian calculations by tracking a lower solution and a penalized
joint solution with first-order updates \citep{kwon2023f2sa}. These methods
differ in utility and approximation, but they all clarify that an attribution
score is tied to an evaluation point or trajectory. Our question is narrower
than a new valuation axiom: under a fixed cardinality budget, does maintaining
and updating the same continuous $\vomega$ improve the subset beyond repeated
local scores at $\vomega^{\mathrm{unif}}$?

\paragraph{Adaptive subset and online batch selection.}
Selection methods modify the data actually consumed during optimization.
Online batch selection prioritizes examples using recently observed losses
\citep{loshchilov2016online}, importance sampling uses inexpensive proxies for
per-example gradient norms \citep{katharopoulos2018importance}, and
GRAD-MATCH periodically chooses subsets whose weighted gradients approximate a
training or validation gradient \citep{killamsetty2021gradmatch}. RHO-LOSS
instead ranks a large candidate batch by current loss minus frozen
irreducible loss and immediately trains on its local top fraction
\citep{mindermann2022rholoss}. Such algorithms couple a scoring rule to a
candidate scope, quota, refresh schedule, and optimizer update, so an end-to-end
comparison cannot identify which choice helps. Our RHO ladder adds global
scoring, batch-local quota, immediate training, online rescoring, and the
frozen comparator separately; the matched batch-size sweep further exposes the
trade-off between local diversity, update granularity, and scoring cost.

\section{Conclusion}

The loss difference shared by ideal RHO-LOSS and the F2SA value-function
construction supplies a precise local bridge between attribution and
selection. The bridge is conditional: it depends on aligned objectives and
adequately tracked inner solutions, and practical RHO-LOSS adds a frozen
comparator approximation as well as a distinct online protocol. The proposed
fixed-horizon experiment makes the decisive comparison explicit: one-shot
selection, repeated reset scoring, and persistent continuous optimization are
evaluated after identical real target training. The complementary RHO ladder
then separates score fidelity from local quota, immediate updates, online
rescoring, and batch scale. Together, these tests can establish not only
whether multi-round selection helps, but whether the gain comes from optimizing
utility beyond the uniform point or from the practical mechanics of online
training.


\bibliography{iclr2027_conference}
\bibliographystyle{iclr2027_conference}

\appendix
\section{Derivation of the Loss-Difference Direction}
\label{app:derivation}

For fixed $\vomega$, let
\begin{equation}
\mathcal{L}^{\alpha,\star}(\vomega)
=\min_{\vtheta}\left\{
f(\vtheta)+\alpha g(\vomega,\vtheta)
\right\}
-\alpha g^\star(\vomega).
\end{equation}
Applying the envelope theorem to the two minimized terms gives
\begin{align}
\nabla_{\vomega}\mathcal{L}^{\alpha,\star}(\vomega)
&=\alpha\nabla_{\vomega}
g(\vomega,\thetatilde(\vomega))
-\alpha\nabla_{\vomega}
g(\vomega,\thetahat(\vomega)).
\end{align}
Since
$g(\vomega,\vtheta)=n^{-1}\sum_i\omega_i\ell(\vtheta;z_i)$,
$\partial g/\partial\omega_i=n^{-1}\ell(\vtheta;z_i)$, which yields
Eq.~(\ref{eq:loss-difference-gradient}). No differentiation through either
argmin is required in this expression; the approximation error instead enters
through finite $\alpha$ and incomplete tracking of the two minimizers.

\section{Interpretation of Reference Points}

The all-present point $\bm{1}$ is natural when asking for the infinitesimal
effect of removing or downweighting an example from the full dataset. The
budget-feasible uniform point $(B/n)\bm{1}$ is natural when asking how to
redistribute a fixed selection budget. A one-shot attribution experiment may
use the former, but a controlled comparison between reset and persistent
budgeted optimization must use the latter. We report both when necessary and
never silently project one into the other.

\section{Go/No-Go Criteria}

The main matrix is launched only if: (i) analytic, finite-difference, unrolled,
and loss-difference gradients agree on the tractable problem; (ii) all weight
states satisfy the capped-budget constraints; (iii) the chosen inner stopping
rule yields stable scores; (iv) no test information is used for stopping or
configuration selection; and (v) every run is reproducible from its saved
configuration, random seeds, and initial subset. If verified VF remains below
practical RHO, we report gradient agreement, inner residuals, score stability,
and constraint audits before interpreting the gap as a method-level result.


\end{document}
